\section{Design Decisions}
\label{sec:design}

\subsection{Text normalisation (shared contract)}
\label{sec:normalisation}

All four implementations apply the same normalisation, in the same order, to
every body:

\begin{enumerate}
    \item lowercase the text;
    \item drop accents/diacritics (Unicode NFD decomposition, combining marks
    removed) --- \texttt{café} $\rightarrow$ \texttt{cafe};
    \item split on any run of non-alphanumeric characters, so punctuation and
    whitespace are separators;
    \item drop tokens that contain a digit (\texttt{chapter3} is discarded);
    \item drop a small stop-word list and tokens shorter than a minimum
    length, and keep the remaining tokens in document order.
\end{enumerate}

A term's postings list is the sorted list of IDs of the documents that
contain it. Serialised, one line of the canonical form is
\texttt{postings\textbackslash ttime\textbackslash t11\textbackslash t84\textbackslash t1342}:
the term \texttt{time} appears in books 11, 84 and 1342.

\subsection{Metadata contract}
\label{sec:metadata-contract}

Headers are parsed with the same rules in the Python pipeline: fields
\texttt{Title:}, \texttt{Author:} and \texttt{Language:} are read from the
Gutenberg header; the language name is mirrored to its ISO~639 code through a
shared table (\texttt{English} $\rightarrow$ \texttt{en}); multiple authors
are split on commas, \texttt{and} and semicolons. The extracted records drive
the SQLite metadata store that the index builders read
(\cref{fig:pipeline}), so indexing goes through the same metadata contract
the production pipeline uses.

The ports mirror the parts of the contract they need: all four split header
and body at the same Gutenberg markers (the Go splitter test pins the same
fixtures as the Python one), and all four apply the identical tokenisation
and postings rules --- which is exactly what the equivalence gate verifies
per language (\cref{sec:equivalence-languages}). Metadata-field extraction
into SQLite is exercised by the Python baseline; the ports consume the
resulting bodies, so the contract they must share is the token/postings one,
and it is byte-verified.

\subsection{Storage choices under evaluation}
\label{sec:storage-choices}

\paragraph{Datalake layouts.}
Time-based layouts favour append-heavy ingestion along a timeline;
book-based layouts favour per-book lookup; batch layouts favour bulk
operations over ID ranges. All three store the same bytes, so the benchmark
can attribute every difference to directory structure alone
(\cref{sec:datalake}).

\paragraph{Index structures.}
The JSON index is a single serialised dictionary (simple, one file, whole
file must be parsed to answer any query). The folder index stores one file
per term (small reads, but one file handle per lookup). MongoDB stores one
document per term with the postings array (server-backed, queried over a
connection). None of these is declared ``best'' a priori: builds, queries,
incremental updates and on-disk footprint are measured for each
(\cref{sec:index}).

\paragraph{Language ports.}
Each port uses its ecosystem's idiomatic tools (for example
\texttt{serde}/\texttt{std::fs} in Rust, \texttt{encoding/json} and
\texttt{os} in Go, \texttt{jakarta.json} plus the driver in Java,
\texttt{json}/\texttt{pathlib}/\texttt{pymongo} in Python). The
\emph{algorithms} are held constant: the same tokenizer semantics, the same
postings construction, the same intersection query, the same synchronous
single-threaded driver. Where a runtime differs by nature (a JIT in Java),
the protocol compensates explicitly (\cref{sec:warmup}); where it cannot
compensate, the difference is reported as a limitation
(\cref{sec:limitations}).
